\chapter{Kết quả}
\label{ch:results}

\section{CIFAR-10 / ResNet-20}

\subsection{Độ chính xác theo mức độ thưa}
\begin{figure}[h]
	\centering
	\includegraphics[width=\textwidth]{resnet20-cifar10_acc+sparsity.png}
	\caption{Độ chính xác theo mức độ thưa trên CIFAR-10/ResNet-20}
	\label{fig:resnet20-cifar10-acc-sparsity}
\end{figure}
Có thể thấy có 2 nhóm hành vi:
 
\paragraph{Nhóm suy giảm từ từ:}

GraSP, SynFlow, Hybrid, Hybrid-Improve và IMP đều duy trì độ chính xác trong khoảng
$85\%$--$91{,}5\%$ qua tất cả các mức độ thưa.
Tại $S=0{,}30$, Hybrid-Improve dẫn đầu ($\approx91{,}2\%$), xấp xỉ bằng mô hình gốc,
trong khi IMP, Hybrid, GraSP và SynFlow dao động trong khoảng $90{,}0\%$--$91{,}0\%$.
Đáng chú ý, sụt giảm của Hybrid-Improve tại $S=0{,}30$ xấp xỉ bằng không, tức là mô
hình sau cắt tỉa nhẹ gần như tương đương mô hình gốc.
Tại $S=0{,}90$, Hybrid và Hybrid-Improve tiếp tục dẫn đầu ($\approx89\%$--$89{,}5\%$),
tiếp theo là GraSP và IMP ($\approx87\%$), trong khi SynFlow sụt giảm đến $\approx85\%$.

\paragraph{Nhóm sụp đổ hiệu năng.}
Early-Bird có hành vi hoàn toàn khác biệt.
Ngay tại $S=0{,}30$, độ chính xác của nó ($\approx89{,}5\%$) đã thấp hơn tất cả các
thuật toán còn lại, và tiếp tục sụt giảm mạnh: $\approx88{,}3\%$ tại $S=0{,}60$, và
$\approx75\%$ tại $S=0{,}90$ --- mất hơn $16$ điểm phần trăm so với mô hình gốc.
Quan trọng hơn, thanh sai số của Early-Bird tại $S=0{,}90$ là rộng nhất trong tất cả các
thuật toán (độ lệch chuẩn $\approx5\%$--$6\%$), cho thấy hành vi của nó ở độ thưa cao
không chỉ kém mà còn không ổn định giữa các hạt giống khác nhau.
Nguyên nhân nằm ở bản chất của Early-Bird: điểm quan trọng kênh được tính từ hệ số
co giãn Batch Normalisation ($\gamma$) tại các epoch rất sớm (thường $6\%$--$20\%$ tổng
số epoch), khi mô hình chưa hội tụ. Tại $S=0{,}90$, các điểm số này quá nhiễu và nhạy
cảm với khởi tạo khác nhau, dẫn đến việc lựa chọn kênh quan trọng sai và không nhất quán.

\subsection{Sụt giảm độ chính xác}

\begin{figure}[h]
	\centering
	\includegraphics[width=\textwidth]{resnet20-cifar10_accdrop.png}
	\caption{Mức giảm độ chính xác so với mô hình dày đặc trên CIFAR-10/ResNet-20}
	\label{fig:resnet20-cifar10-accdrop}
\end{figure}
\label{subsec:r20c10_drop}
 
Hình~\ref{fig:resnet20-cifar10-accdrop} biểu diễn sụt giảm độ chính xác trung bình và độ lệch chuẩn
qua các hạt giống cho mỗi thuật toán tại từng mức độ thưa.
 
Tại $S=0{,}30$, sụt giảm của tất cả các thuật toán ngoài Early-Bird đều nhỏ và tập trung:
Hybrid-Improve $\approx0{,}0\%$, IMP $\approx0{,}5\%$, Hybrid $\approx0{,}9\%$,
SynFlow $\approx1{,}4\%$ và GraSP $\approx1{,}5\%$.
Thanh lỗi nhỏ ở mức này xác nhận rằng $S=0{,}30$ nằm trong vùng nén an toàn cho
ResNet-20 với tất cả các phương pháp không có cấu trúc.
 
Tại $S=0{,}60$, thứ hạng được giữ nguyên: Hybrid-Improve vẫn dẫn đầu ($\approx0{,}5\%$),
tiếp theo là Hybrid ($\approx1{,}1\%$), IMP ($\approx1{,}3\%$), SynFlow ($\approx1{,}7\%$)
và GraSP ($\approx2{,}4\%$).
 
Điều quan trọng nhất nằm ở $S=0{,}90$.
Early-Bird đạt sụt giảm trung bình $\approx17\%$ với độ lệch chuẩn $\approx6\%$ --- một
mức biến động cực kỳ cao, có nghĩa là trong điều kiện tệ nhất, mô hình có thể mất hơn
$23\%$ độ chính xác.
Trong số các phương pháp còn lại, Hybrid dẫn đầu với chỉ $\approx2{,}2\%$ sụt giảm,
tiếp theo là Hybrid-Improve ($\approx3{,}0\%$), IMP ($\approx5{,}0\%$), GraSP
($\approx5{,}5\%$) và SynFlow ($\approx6{,}2\%$).
 
Đáng chú ý, Hybrid ở $S=0{,}90$ vượt qua Hybrid-Improve theo
giá trị trung bình. Tuy nhiên, biên độ độ chính xác của Hybrid lớn hơn đáng kể so với 
Hybrid-Improve, do đó vẫn có thể nói rằng Hybrid-Improve có hiệu suất ổn định hơn ở mức độ thưa cực cao.
\subsection{Mức giảm FLOPs vs chi phí huấn luyện và độ chính xác}
\label{subsec:r20c10_efficiency}
\noindent
\begin{figure}[h]
	\centering
	\includegraphics[width=\textwidth]{resnet20-cifar10_pareto.png}
	\caption{Biểu đồ Pareto: Cân bằng giữa độ chính xác và độ trễ trên CIFAR-10/ResNet-20}
	\label{fig:resnet20-cifar10-pareto}
\end{figure}

Hình~\ref{fig:resnet20-cifar10-pareto} trình bày hai biểu đồ phân tán: (trái) độ chính xác so
với mức giảm FLOPs tiết kiệm, và (phải) chi phí huấn luyện so với mức giảm FLOPs.
Cả hai đều hiển thị giá trị trung bình và thanh lỗi qua nhiều hạt giống.
 
Cần lưu ý rằng mức giảm FLOPs được ước lượng theo mật độ mặt nạ và là các giá trị lý thuyết, không phải phép đo phần
cứng thực tế.
 
Tại tỉ lệ giảm FLOPs $\approx20\%$--$50\%$, GraSP, Hybrid và Hybrid-Improve
tập trung gần phía trên của trục độ chính xác ($\approx90\%$--$91\%$), cho thấy chúng
là các phương pháp hiệu quả nhất khi nén nhẹ.
Ở tỉ lệ giảm FLOPs $>80\%$, Hybrid và Hybrid-Improve
nằm trên biên Pareto với $\approx89\%$--$89{,}5\%$, tiếp theo là GraSP và IMP ở
$\approx87\%$.
 
Về chi phí huấn luyện, IMP là đắt nhất ở mọi mức độ thưa ($10^4$--$10^5$ giây), do có
nhiều vòng prune-retrain-rewind. Hybrid và Hybrid-Improve ở mức trung gian, trong khi
GraSP và SynFlow rẻ nhất vì chỉ cần 1 lần cắt tỉa và một lần huấn luyện.
Tuy nhiên, kết quả cho thấy IMP có biến động chi phí thấp, trong
khi Hybrid và Hybrid-Improve có biến động vừa phải tùy thuộc vào khi nào điều kiện dừng
sớm (early stopping) được kích hoạt. % original hybrid paper suggests this

\subsection{Phân phối độ thưa theo tầng và tính ổn định}
\label{subsec:r20c10_layerwise}

\begin{figure}[h]
    \centering
    \begin{subfigure}{0.49\textwidth}
        \centering
        \includegraphics[width=\textwidth]{resnet20-cifar10_s0.3_heatmap.png}
        \caption{Bản đồ nhiệt phân phối độ thưa theo tầng tại $S=0{,}30$}
        \label{fig:r20c10_heatmap_s03}
    \end{subfigure}
    \hfill
    \begin{subfigure}{0.49\textwidth}
        \centering
        \includegraphics[width=\textwidth]{resnet20-cifar10_s0.6_heatmap.png}
        \caption{Bản đồ nhiệt phân phối độ thưa theo tầng tại $S=0{,}60$}
        \label{fig:r20c10_heatmap_s06}
    \end{subfigure}
    \hfill
    \begin{subfigure}{0.49\textwidth}
        \centering
        \includegraphics[width=\textwidth]{resnet20-cifar10_s0.9_heatmap.png}
        \caption{Bản đồ nhiệt phân phối độ thưa theo tầng tại $S=0{,}90$}
        \label{fig:r20c10_heatmap_s09}
    \end{subfigure}

    \caption{Bản đồ nhiệt phân phối độ thưa theo tầng trên CIFAR-10/ResNet-20}
    \label{fig:grouped-accuracy}
\end{figure}
 
Hình~\ref{fig:r20c10_heatmap_s03}, \ref{fig:r20c10_heatmap_s06} và
\ref{fig:r20c10_heatmap_s09} hiển thị cả bản đồ nhiệt trung bình và
độ lệch chuẩn của phân phối độ thưa theo tầng qua các cài đặt.
 
GraSP cho thấy phân phối độ thưa đồng đều nhất qua tất cả các tầng ở cả ba mức độ thưa,
với độ lệch chuẩn rất thấp trên hầu hết các tầng.
Tính đồng đều trong phân phối của GraSP là hệ quả trực tiếp từ tiêu chí bảo toàn chuẩn gradient: vì chuẩn gradient
phụ thuộc vào luồng thông tin qua toàn bộ mạng, bất kỳ tầng nào bị cắt tỉa quá mạnh đều tạo ra
nút thắt cổ chai làm giảm gradient toàn cục, khiến các trọng số tại tầng đó nhận điểm saliency
cao và được bảo vệ, điều này tự nhiên khuyến khích cắt tỉa đồng đều trên tất cả các tầng.
 
IMP cũng có độ lệch chuẩn rất thấp qua các tầng, điều này có thể giải thích bởi tính xác định
cao của tiêu chí magnitude sau huấn luyện đầy đủ — các trọng số hội tụ về vùng tương tự qua
các hạt giống, dẫn đến ngưỡng cắt tỉa toàn cục nhất quán. % Morcos et al. (2019) "One ticket to win them all" cho thấy winning ticket có thể transfer qua datasets
Tuy nhiên, IMP tập trung độ thưa cao hơn vào các tầng đầu (ví dụ
\texttt{layer1.0.conv1}: $55{,}1\%$ tại $S=0{,}30$), Điều này có thể giải thích do
các tầng đầu của ResNet-20 trên CIFAR-10 thường học các đặc trưng mức thấp như cạnh
và kết cấu đơn giản, do đó có thể chứa nhiều trọng số dư thừa hoặc có độ lớn nhỏ. Vì
vậy, các tầng này thường bị cắt tỉa mạnh hơn.
Cả hai biến thể hybrid đều thể hiện xu hướng bảo vệ tầng đầu: tầng \texttt{conv1}
chỉ đạt $\approx13\%$--$17\%$ độ thưa tại $S=0{,}30$, trong khi các tầng sâu hơn nhận
được cắt tỉa nhiều hơn (ví dụ \texttt{layer3.2.conv2}: $55\%$--$59{,}5\%$).
Độ lệch chuẩn ở mức trung bình --- 
phản ánh sự kết hợp của hai giai đoạn cắt tỉa (one-shot và lặp). 

SynFlow là thuật toán kém ổn định nhất trong số các phương pháp được đánh giá.
Tại $S=0{,}30$, độ lệch chuẩn đạt tới $52{,}6\%$ tại \texttt{layer1.0.conv1} và $28{,}7\%$
tại \texttt{conv1}. Điều này có nghĩa là quyết định phân bổ độ thưa cho các tầng quan
trọng này thay đổi rất lớn giữa các hạt giống.
Nguyên nhân nằm ở tính nhạy cảm của điểm SynFlow đối với giá trị khởi tạo ngẫu nhiên:
vì điểm path-norm được tính tại thời điểm khởi tạo mà không dùng dữ liệu thực, các trọng % 
số được khởi tạo khác nhau sẽ tạo ra phân phối điểm rất khác nhau.
Sự không ổn định trong phân phối độ thưa theo tầng trực tiếp giải thích sự biến động cao
trong độ chính xác cuối cùng của SynFlow qua các hạt giống.
 
Bản đồ độ lệch chuẩn của Early-Bird cho thấy biến động lớn ở nhiều tầng, đặc biệt ở các
kết nối shortcut và các tầng conv đầu (ví dụ \texttt{layer3.0.shortcut.0}: $\sigma=26{,}6\%$ % wrong data
tại $S=0{,}30$).
Điều này xác nhận rằng các điểm quan trọng kênh dựa trên hệ số $\gamma$ BN ở giai đoạn
huấn luyện sớm không ổn định và phụ thuộc nhiều vào hạt giống khởi tạo, dẫn đến các quyết
định cắt tỉa cấu trúc khác nhau và do đó độ chính xác khác nhau đáng kể giữa các lần chạy


\section{CIFAR-100/ResNet-20}

CIFAR-100 là bộ dữ liệu khó hơn đáng kể so với CIFAR-10: $100$ lớp thay vì $10$ lớp,
với $600$ ảnh mỗi lớp thay vì $6000$, đòi hỏi mô hình phải học các đặc trưng phân biệt
tinh tế hơn với ít dữ liệu huấn luyện hơn mỗi lớp.
Khi áp dụng cùng các thuật toán cắt tỉa trên cấu hình ResNet-20/CIFAR-100, tác động của cắt tỉa sẽ nghiêm trọng hơn, đặc biệt ở mức độ thưa cao, vì mạng cần
duy trì khả năng phân biệt giữa nhiều lớp hơn với ít tham số hơn.
 
\subsection{Độ chính xác và sụt giảm theo mức độ thưa}
\label{subsec:r20c100_acc}

\begin{figure}[h]
    \centering

    \begin{subfigure}{0.49\textwidth}
        \centering
        \includegraphics[width=\textwidth]{resnet20-cifar100_acc+sparsity.png}
        \caption{Độ chính xác theo mức độ thưa trên CIFAR-100/ResNet-20}
        \label{fig:resnet20-cifar100-acc-sparsity}
    \end{subfigure}
    \hfill
    \begin{subfigure}{0.49\textwidth}
        \centering
        \includegraphics[width=\textwidth]{resnet20-cifar100_accdrop.png}
        \caption{Mức giảm độ chính xác so với mô hình dày đặc trên CIFAR-100/ResNet-20}
        \label{fig:resnet20-cifar100-accdrop}
    \end{subfigure}

    \caption{Độ chính xác theo mức độ thưa trên CIFAR-100/ResNet-20}
    \label{fig:grouped-accuracy}
\end{figure}
 
Hình~\ref{fig:resnet20-cifar100-acc-sparsity} và \ref{fig:resnet20-cifar100-accdrop} cho thấy xu hướng tương tự như
CIFAR-10 nhưng với mức sụt giảm lớn hơn đáng kể.
 
Tại $S=0{,}30$, Hybrid và Hybrid-Improve gần như không sụt giảm ($\approx1{,}0\%$ và
$\approx0{,}7\%$ tương ứng), trong khi IMP sụt giảm $\approx0{,}7\%$, SynFlow $\approx2{,}5\%$,
GraSP và Early-Bird cùng $\approx3{,}8\%$.
Thanh lỗi nhỏ cho hầu hết các thuật toán ở mức này cho thấy $S=0{,}30$ vẫn trong vùng
nén an toàn trên CIFAR-100 với các phương pháp không có cấu trúc.
 
Tại $S=0{,}60$, khoảng cách giữa các thuật toán bắt đầu mở rộng rõ rệt.
Early-Bird sụt giảm $\approx8{,}5\%$ với thanh lỗi lớn, cho thấy độ không ổn định đặc
trưng của nó.
Hybrid ($\approx2{,}2\%$) và Hybrid-Improve ($\approx2{,}5\%$) vẫn duy trì hiệu suất tốt.
GraSP ($\approx4\%$), IMP ($\approx3\%$) và SynFlow ($\approx4\%$) ở mức trung gian.
 
Tại $S=0{,}90$, sự phân hóa trở nên cực kỳ rõ ràng:
\begin{itemize}
    \item \textbf{Early-Bird:} sụt giảm trung bình $\approx35\%$ với độ lệch chuẩn lên
    đến $\approx13\%$. Trong điều kiện tệ nhất (mean + 1 std), độ chính xác có thể rơi
    xuống dưới $20\%$ --- gần như không thể sử dụng được. Đây là hành vi sụp đổ nghiêm
    trọng nhất trong toàn bộ thực nghiệm.
    \item \textbf{GraSP và IMP:} cùng sụt giảm $\approx13\%$ --- lớn hơn gấp đôi so với
    CIFAR-10 ở cùng mức độ thưa.
    \item \textbf{SynFlow:} sụt giảm $\approx14{,}5\%$ --- lớn hơn đáng kể so với CIFAR-10
    ($\approx6{,}2\%$).
    \item \textbf{Hybrid và Hybrid-Improve:} cùng sụt giảm $\approx6\%$, dẫn đầu so với các
    thuật toán khác với biên độ tương đối nhỏ.
\end{itemize}
 
Việc mức sụt giảm lớn hơn đối với GraSP, IMP và SynFlow trên CIFAR-100 là có thể dự đoán được.
Trên CIFAR-100, mỗi tham số mang nhiều thông tin phân loại hơn
do phải phân biệt $10\times$ nhiều lớp hơn với $10\times$ ít dữ liệu hơn mỗi lớp hơn so với CIFAR-10.
Cắt tỉa $90\%$ tham số do đó loại bỏ nhiều thông tin không thể phục hồi hơn, và khả năng
phục hồi thông qua tinh chỉnh bị hạn chế hơn vì tập dữ liệu nhỏ hơn theo lớp.
Các biến thể hybrid vẫn vượt trội nhờ chiến lược phân giai đoạn: cắt tỉa từ từ và tinh
chỉnh với kiên nhẫn sau mỗi bước đặc biệt quan trọng khi tập dữ liệu nhỏ hơn theo lớp,
vì mô hình cần nhiều epoch hơn để hội tụ lại sau mỗi bước cắt tỉa.
 
\subsection{Phân phối độ thưa theo tầng và tính ổn định}
\label{subsec:r20c100_layerwise}

\begin{figure}[h]
	\centering
	\begin{subfigure}{0.49\textwidth}
		\centering
		\includegraphics[width=\textwidth]{resnet20-cifar100_s0.3_heatmap.png}
		\caption{Bản đồ nhiệt phân phối độ thưa theo tầng tại $S=0{,}30$}
		\label{fig:r20c100_heatmap_s03}
	\end{subfigure}
	\hfill
	\begin{subfigure}{0.49\textwidth}
		\centering
		\includegraphics[width=\textwidth]{resnet20-cifar100_s0.6_heatmap.png}
		\caption{Bản đồ nhiệt phân phối độ thưa theo tầng tại $S=0{,}60$}
		\label{fig:r20c100_heatmap_s06}
	\end{subfigure}
	\hfill
	\begin{subfigure}{0.49\textwidth}
		\centering
		\includegraphics[width=\textwidth]{resnet20-cifar100_s0.9_heatmap.png}
		\caption{Bản đồ nhiệt phân phối độ thưa theo tầng tại $S=0{,}90$}
		\label{fig:r20c100_heatmap_s09}
	\end{subfigure}

	\caption{Bản đồ nhiệt phân phối độ thưa theo tầng trên CIFAR-100/ResNet-20}
	\label{fig:grouped-accuracy}
\end{figure}
 
Hình~\ref{fig:r20c100_heatmap_s03}, \ref{fig:r20c100_heatmap_s06} và
\ref{fig:r20c100_heatmap_s09} cho thấy các xu hướng phân phối tương tự như CIFAR-10 nhưng
với một số điểm khác biệt đáng chú ý.

Early-Bird trên CIFAR-100 tại $S=0{,}30$ đã có nhiều tầng đạt $>90\%$ độ thưa kênh
(ví dụ \texttt{layer1.2.conv2}: $96{,}1\%$; \texttt{layer3.0.shortcut.0}: $95{,}8\%$),
sớm và cực đoan hơn nhiều so với CIFAR-10.
Điều này cho thấy điểm $\gamma$ BN trên CIFAR-100 tạo ra phân phối thưa kênh nghiêm trọng
hơn ngay cả ở mức nén nhẹ --- một phần vì mô hình trên CIFAR-100 học được ít kênh ``dư
thừa'' hơn, khiến các điểm $\gamma$ phân biệt rõ hơn và dẫn đến cắt tỉa hung hăng hơn.
 
SynFlow trên CIFAR-100 duy trì mô hình độ lệch chuẩn cao ở các tầng đầu
(\texttt{conv1}: $\sigma=26{,}5\%$; \texttt{layer1.0.conv1}: $\sigma=52{,}0\%$ tại $S=0{,}30$),
xác nhận rằng sự không ổn định của SynFlow là đặc trưng của phương pháp chứ không phải
của bộ dữ liệu cụ thể.
 
\subsection{Mức giảm FLOPs vs chi phí huấn luyện và độ chính xác}
\label{subsec:r20c100_efficiency}

\begin{figure}[h]
	\centering
	\includegraphics[width=\textwidth]{resnet20-cifar100_pareto.png}
	\caption{Biểu đồ Pareto: Cân bằng giữa độ chính xác và độ trễ trên CIFAR-100/ResNet-20}
	\label{fig:resnet20-cifar100-pareto}
\end{figure}
 
Hình~\ref{fig:resnet20-cifar100-pareto} cho thấy bức tranh hiệu quả tương tự như CIFAR-10 nhưng
với dải độ chính xác thấp hơn ($35\%$--$67\%$ thay vì $75\%$--$91\%$).
 
Tại tỉ lệ giảm FLOPs cao ($>80\%$), Hybrid và Hybrid-Improve đạt $\approx61{,}5\%$ độ chính
xác với chi phí huấn luyện trung bình, trong khi GraSP đạt $\approx55\%$ với chi phí thấp
hơn đáng kể.
IMP đạt $\approx54{,}5\%$ với chi phí rất cao --- càng khẳng định rằng chi phí
cao của IMP không tương xứng với hiệu suất ở mức độ thưa cực cao.
SynFlow đạt $\approx53\%$ với chi phí thấp, nhưng với độ biến động cao hơn.

\section{CIFAR-10/ResNet-50}

\subsection{Độ chính xác và sụt giảm theo mức độ thưa}
\label{subsec:r50_acc}
\begin{figure}
	\centering
	\begin{subfigure}{0.49\textwidth}
		\centering
		\includegraphics[width=\textwidth]{resnet50-cifar10_acc+sparsity.png}
		\caption{Độ chính xác theo mức độ thưa trên CIFAR-10/ResNet-50}
		\label{fig:resnet50-cifar10-acc-sparsity}
	\end{subfigure}
	\hfill
	\begin{subfigure}{0.49\textwidth}
		\centering
		\includegraphics[width=\textwidth]{resnet50-cifar10_accdrop.png}
		\caption{Mức giảm độ chính xác so với mô hình dày đặc trên CIFAR-10/ResNet-50}
		\label{fig:resnet50-cifar10-accdrop}
	\end{subfigure}

	\caption{Độ chính xác theo mức độ thưa trên CIFAR-10/ResNet-50}
	\label{fig:grouped-accuracy}
\end{figure}
 
Hình~\ref{fig:resnet50-cifar10-acc-sparsity} và \ref{fig:resnet50-cifar10-accdrop} cho thấy các hiện tượng sau:

Sụt giảm âm $\approx-0{,}3\%$ của Hybrid-Improve tại $S=0{,}30$ không phải là sai số
đo lường --- đây là hiệu ứng điều chỉnh ngầm trong vùng dư thừa tham số, nơi cắt tỉa
nhẹ loại bỏ nhiễu và cải thiện khả năng tổng quát hóa.
Hybrid cũng đạt sụt giảm thấp ($\approx1{,}4\%$).
 
Hybrid-Improve tiếp tục dẫn đầu ($\approx1\%$ sụt giảm), tiếp theo là Hybrid
($\approx2\%$), GraSP ($\approx2\%$), Early-Bird ($\approx3\%$), IMP ($\approx3{,}5\%$)
và SynFlow ($\approx3{,}5\%$).
Sự nhất quán qua các kiến trúc củng cố kết luận rằng chiến lược hybrid mạnh mẽ ở mức nén
vừa phải bất kể độ sâu mạng.
 
Sụt giảm của IMP tăng từ $\approx5\%$ trên ResNet-20 lên $\approx7{,}5\%$ trên ResNet-50,
trong khi cấu hình ($10$ vòng lặp, $160$ epoch/vòng) không được tinh chỉnh lại.
Trên ResNet-50 ($\approx25$M tham số), $10$ vòng lặp phải xử lý không gian trọng số phức
tạp hơn nhiều, và điều kiện ổn định cần thiết cho cơ chế rewind khó đạt được hơn trong
cùng số epoch.
 
Sụt giảm của SynFlow tăng từ $\approx6{,}2\%$ trên ResNet-20 lên $\approx10{,}5\%$ ---
mức tăng lớn nhất trong tất cả các thuật toán không phải Early-Bird.
Như được ghi nhận trong bản đồ nhiệt (Hình~\ref{fig:r50_heatmap_s09}), tại $S=0{,}90$
nhiều tầng shortcut đạt $\approx100\%$ độ thưa, loại bỏ hiệu quả các kết nối residual và
biến ResNet-50 thành một mạng phẳng bị hư hại nghiêm trọng.
ResNet-50 có nhiều bottleneck block hơn gấp ba lần ResNet-20, mỗi block có một shortcut
chiếu $1{\times}1$, làm trầm trọng thêm xu hướng của SynFlow khi tập trung độ thưa vào các
kết nối này.
 
Hybrid đạt $\approx3{,}3\%$ sụt giảm độ chính xác tại $S=0{,}90$ trên ResNet-50 --- tốt hơn so với
ResNet-20.
Hybrid-Improve đạt $\approx3{,}5\%$ --- hầu như không đổi so với ResNet-20.
Sự ổn định giữa các kiến trúc này là kết quả quan trọng: chiến lược cắt tỉa của 2 thuật toán này không  
có bước reset weight như IMP (cũng sử dụng magnitude như tiêu chí cắt tỉa) cho phép mạng thích nghi dần
thay vì phải học lại từ đầu sau mỗi vòng như IMP.

Early-Bird sụt giảm ~19\% trên ResNet-50, tệ hơn so với ~17\% trên ResNet-20,
mặc dù ResNet-50 có kênh rộng hơn cung cấp nhiều
dư địa hơn trước. Hành vi sụp đổ vẫn xảy ra do điểm $\gamma$ BN tại giai đoạn sớm 
không đủ tin cậy ở $S=0{,}90$.
 
\subsection{Phân phối độ thưa theo tầng và tính ổn định}
\label{subsec:r50_layerwise}

\begin{figure}[h]
	\centering
	\begin{subfigure}{0.49\textwidth}
		\centering
		\includegraphics[width=\textwidth]{resnet50-cifar10_s0.3_heatmap.png}
		\caption{Bản đồ nhiệt phân phối độ thưa theo tầng tại $S=0{,}30$}
		\label{fig:r50_heatmap_s03}
	\end{subfigure}
	\hfill
	\begin{subfigure}{0.49\textwidth}
		\centering
		\includegraphics[width=\textwidth]{resnet50-cifar10_s0.6_heatmap.png}
		\caption{Bản đồ nhiệt phân phối độ thưa theo tầng tại $S=0{,}60$}
		\label{fig:r50_heatmap_s06}
	\end{subfigure}
	\hfill
	\begin{subfigure}{0.49\textwidth}
		\centering
		\includegraphics[width=\textwidth]{resnet50-cifar10_s0.9_heatmap.png}
		\caption{Bản đồ nhiệt phân phối độ thưa theo tầng tại $S=0{,}90$}
		\label{fig:r50_heatmap_s09}
	\end{subfigure}

	\caption{Bản đồ nhiệt phân phối độ thưa theo tầng trên CIFAR-10/ResNet-50}
	\label{fig:grouped-accuracy}
\end{figure}
 
Hình~\ref{fig:r50_heatmap_s03}, \ref{fig:r50_heatmap_s06} và \ref{fig:r50_heatmap_s09}
cho thấy các mô hình định tính tương tự như ResNet-20 nhưng với sự khác biệt về cường độ.
 
IMP trên ResNet-50 tại $S=0{,}30$ đã đẩy một số tầng shortcut đầu lên gần $100\%$ độ thưa
ngay cả ở $S=0{,}30$ --- cực đoan hơn nhiều so với ResNet-20 ở cùng mức nén.
Trong kiến trúc bottleneck, các tầng này chủ yếu đảm nhiệm vai trò mở rộng số chiều kênh và 
biến đổi biểu diễn đặc trưng. Trên CIFAR-10, mức độ mở rộng này có thể dư thừa so với độ phức 
tạp của bài toán, dẫn đến nhiều trọng số có độ lớn nhỏ và dễ bị loại bỏ.
 
GraSP tiếp tục là phương pháp phân phối đồng đều nhất, duy trì mức độ thưa tương đương
qua tất cả các tầng ngay cả trong kiến trúc sâu hơn, phản ánh rằng tiêu chí bảo toàn
gradient tìm kiếm sự đồng đều về đóng góp saliency mà không bị ảnh hưởng bởi sự khác
biệt về kích thước tầng.
 
SynFlow tại $S=0{,}90$ cho thấy nhiều tầng shortcut đạt $\approx99\%$--$100\%$ độ thưa,
xác nhận giả thuyết sụp đổ cấu trúc tại kết nối residual.


\subsection{Mức giảm FLOPs vs chi phí huấn luyện và độ chính xác}
Kết quả trên ResNet-50 có một số điểm khác biệt quan trọng so
với ResNet-20.

\begin{figure}
	\centering
	\includegraphics[width=\textwidth]{resnet50-cifar10_pareto.png}
	\caption{Biểu đồ Pareto: Cân bằng giữa độ chính xác và độ trễ trên CIFAR-10/ResNet-50}
	\label{fig:resnet50-cifar10-pareto}
\end{figure}

Tại $S=0{,}90$ ($\approx85\%$--$90\%$ FLOPs tiết kiệm), GraSP đạt $\approx89\%$
độ chính xác --- xấp xỉ bằng Hybrid ($\approx89{,}7\%$) và Hybrid-Improve
($\approx89{,}2\%$) --- nhưng với chi phí huấn luyện thấp hơn khoảng một bậc
độ lớn.
Trên ResNet-20, khoảng cách accuracy giữa GraSP và các biến thể hybrid tại
$S=0{,}90$ là $\approx2$--$3$ điểm phần trăm, đủ để biện minh cho chi phí cao
hơn của hybrid.
Trên ResNet-50, khoảng cách này thu hẹp xuống dưới $1$ điểm phần trăm, trong
khi chi phí huấn luyện vẫn chênh lệch lớn.
Điều này có nghĩa là trên ResNet-50/CIFAR-10, GraSP chiếm \emph{vị trí tốt hơn
trên biên Pareto} so với Hybrid và Hybrid-Improve khi xét đồng thời cả accuracy
lẫn chi phí huấn luyện.

Lý giải cho sự cải thiện tương đối này của GraSP trên ResNet-50 nằm ở đặc tính
bảo toàn chuẩn gradient của nó: với mạng sâu hơn và nhiều tham số hơn, tiêu chí
Hessian-gradient product của GraSP tạo ra phân phối độ thưa đồng đều hơn giữa
các tầng, tránh được các nút thắt cổ chai mà magnitude pruning toàn cục (IMP,
Hybrid) dễ tạo ra khi không gian trọng số lớn hơn nhiều \cite{grasp}.
 
% ===========================================================================
\section{Độ trễ suy diễn}
\label{sec:latency}
% ===========================================================================
\begin{figure}
	\centering
	\begin{subfigure}{0.49\textwidth}
		\centering
		\includegraphics[width=\textwidth]{resnet20-cifar10_latency.png}
		\caption{Độ trễ suy diễn trên CIFAR-10/ResNet-20}
		\label{fig:r20c10_latency}
	\end{subfigure}
	\hfill
	\begin{subfigure}{0.49\textwidth}
		\centering
		\includegraphics[width=\textwidth]{resnet20-cifar100_latency.png}
		\caption{Độ trễ suy diễn trên CIFAR-100/ResNet-20}
		\label{fig:r20c100_latency}
	\end{subfigure}
	\hfill
	\begin{subfigure}{0.49\textwidth}
		\centering
		\includegraphics[width=\textwidth]{resnet50-cifar10_latency.png}
		\caption{Độ trễ suy diễn trên CIFAR-10/ResNet-50}
		\label{fig:r50_latency}
	\end{subfigure}

	\caption{So sánh độ trễ suy diễn giữa mô hình gốc và mô hình đã cắt tỉa trên các cấu hình khác nhau}
	\label{fig:grouped-accuracy}
\end{figure}
 
Hình~\ref{fig:r20c10_latency}, \ref{fig:r20c100_latency} và \ref{fig:r50_latency} trình
bày so sánh độ trễ suy diễn (ms/batch) giữa mô hình gốc (hình tròn) và mô hình đã cắt tỉa
(hình vuông) cho mỗi thuật toán và mức độ thưa.
Quan sát nổi bật nhất là tất cả các
phương pháp không có cấu trúc đều có thanh lỗi rất lớn về độ trễ, đặc biệt ở mức độ thưa
thấp và vừa.
Tại $S=0{,}30$ và $S=0{,}60$ trên ResNet-20, thanh lỗi của nhiều thuật toán (SynFlow,
Hybrid, Hybrid-Improve, GraSP, Early-Bird) trải rộng hơn $2$ ms --- gần bằng toàn bộ dải
giá trị đo được.
Điều này có nghĩa là sự khác biệt quan sát được giữa mô hình gốc và mô hình đã cắt tỉa
nằm hoàn toàn trong vùng nhiễu đo lường.
Đến $S=0{,}90$, các thanh lỗi thu hẹp đáng kể (khoảng $\pm0{,}1$--$0{,}3$ ms), nhưng
lúc này sự khác biệt giữa dense và pruned cũng biến mất --- mọi thuật toán tập trung trong
khoảng $1{,}8$--$2{,}6$ ms/batch với ResNet-20.


